%% kapa: exam=1 round=25 year=2014 period=1 form=A subject=civil_law subject_ko=민법 seq=12 number=12
%% source: Q-Net 감정평가사 기출문제 (cst00309) artlSeq=5101207 file=extracted/1008/5101207/2014년 제25회 감정평가사 1차시험 기출문제/1교시(A형).hwp
%% conversion: kordoc 4.24.51 -> tools/gen1.py (auto); review=proofread-sonnet; flags=none
\begin{question}{12}
점유보호청구권에 관한 설명으로 옳은 것은? (다툼이 있으면 판례에 의함)
\begin{choices}
  \item 직접점유자가 제3자로부터 그 점유를 침탈당하더라도 간접점유자는 점유보호청구권을 행사할 수 없다.
  \item 유실물을 우연히 습득한 자에 대해서는 점유물반환청구권을 행사할 수 없다.
  \item 점유물반환청구의 상대방은 현재 점유를 침해하고 있는 자이므로, 점유침탈자의 선의의 특별승계인도 그 상대방이 된다.
  \item 영업상 타인의 지시를 받아 물건을 사실상 지배하는 자에게도 점유보호청구권이 인정된다.
  \item 점유물반환청구권의 행사기간은 출소기간이 아니다.
\end{choices}
\end{question}
\answer{2}
